\documentclass[11pt,a4paper]{article}
\usepackage[utf8]{inputenc}
\usepackage[T1]{fontenc}
\usepackage[french]{babel}
\usepackage{lmodern}
\usepackage{amsmath,amssymb}
\usepackage[margin=2.2cm]{geometry}
\usepackage{enumitem}
\usepackage{tkz-tab}
\usepackage{listings}
\usepackage{fancyhdr}

\lstset{
  language=Python,
  basicstyle=\ttfamily\small,
  frame=single,
  columns=fullflexible,
  keepspaces=true,
  showstringspaces=false
}

\pagestyle{fancy}
\fancyhf{}
\lhead{Corrigé -- Fiche d'exercices pour Maximilien (8)}
\rhead{\thepage}

\newcommand{\e}{\mathrm{e}}
\newcommand{\R}{\mathbb{R}}
\setlist[enumerate,1]{label=\arabic*.}
\setlist[enumerate,2]{label=(\alph*)}

\newcommand{\exo}[1]{\section*{#1}}
\newcommand{\rep}{\par\smallskip\noindent\textbf{Réponse. }}

\begin{document}

\begin{center}
{\Large\bfseries Corrigé détaillé de la fiche d'exercices (8)}\\[2pt]
28/10/2026
\end{center}

%==========================================================
\exo{Exercice 1 (Métropole 1 2026)}
%==========================================================

On a $u_0 = 20$ et, pour tout $n\in\mathbb{N}$,
$u_{n+1} = 0{,}965\,u_n + 0{,}35 + 0{,}07\,\e^{-0{,}1n}$.

\begin{enumerate}

\item \textbf{Montrons que $u_1 = 19{,}72$.}

On applique la relation de récurrence avec $n=0$ :
\[
u_1 = 0{,}965\times u_0 + 0{,}35 + 0{,}07\,\e^{0}
= 0{,}965\times 20 + 0{,}35 + 0{,}07\times 1
= 19{,}3 + 0{,}35 + 0{,}07 = 19{,}72.
\]

\item \textbf{Montrons par récurrence que, pour tout $n\in\mathbb{N}$, $u_n > 10$.}

Soit $P(n)$ la propriété : « $u_n > 10$ ».

\emph{Initialisation.} $u_0 = 20 > 10$, donc $P(0)$ est vraie.

\emph{Hérédité.} Soit $n\in\mathbb{N}$ tel que $P(n)$ soit vraie, c'est-à-dire $u_n > 10$.
Alors $0{,}965\,u_n > 0{,}965\times 10 = 9{,}65$, d'où
\[
u_{n+1} = 0{,}965\,u_n + 0{,}35 + 0{,}07\,\e^{-0{,}1n}
> 9{,}65 + 0{,}35 + 0{,}07\,\e^{-0{,}1n}
= 10 + 0{,}07\,\e^{-0{,}1n}.
\]
Or $0{,}07\,\e^{-0{,}1n} > 0$ car une exponentielle est strictement positive, donc $u_{n+1} > 10$ : $P(n+1)$ est vraie.

\emph{Conclusion.} Par récurrence, pour tout $n\in\mathbb{N}$, $u_n > 10$.

\item \textbf{Convergence de $(u_n)$.}

La suite $(u_n)$ est décroissante (admis) et minorée par $10$ (question 2).
D'après le théorème de convergence monotone, \emph{toute suite décroissante et minorée converge}.
Donc $(u_n)$ est convergente.

\item \textbf{Limite de $(u_n)$.}

\begin{enumerate}
\item Notons $\ell$ la limite de $(u_n)$. Alors $\lim_{n\to+\infty}u_{n+1}=\ell$ également.
De plus, $\lim_{n\to+\infty}\e^{-0{,}1n}=0$ car $-0{,}1n\to-\infty$.
Par opérations sur les limites, le membre de droite de $u_{n+1} = 0{,}965\,u_n + 0{,}35 + 0{,}07\,\e^{-0{,}1n}$
tend vers $0{,}965\,\ell + 0{,}35$. Par unicité de la limite :
\[
\ell = 0{,}965\,\ell + 0{,}35.
\]
Donc $\ell$ est solution de l'équation $x = 0{,}965\,x + 0{,}35$.

\item On résout :
\[
x = 0{,}965\,x + 0{,}35 \iff 0{,}035\,x = 0{,}35 \iff x = \dfrac{0{,}35}{0{,}035} = 10.
\]
Donc $\ell = 10$.

\emph{Interprétation.} À long terme, la température de la pièce se stabilise à $10\,^{\circ}\mathrm{C}$
(température vers laquelle la pièce tend à se refroidir sans chauffage).
\end{enumerate}

\item \textbf{Remise en marche du chauffage.}

\begin{enumerate}
\item Le chauffage se remet en marche dès que $u_n \leqslant 18$.
On boucle donc tant que $u > 18$ (la fonction \texttt{exp} provient du module \texttt{math}) :
\begin{lstlisting}
from math import exp

def marche():
    n = 0
    u = 20
    while u > 18:
        u = 0.965*u + 0.35 + 0.07*exp(-0.1*n)
        n = n + 1
    return n
\end{lstlisting}
Ligne 4 : \texttt{while u > 18 :}\quad
Ligne 5 : \texttt{u = 0.965*u + 0.35 + 0.07*exp(-0.1*n)}\quad
Ligne 6 : \texttt{n = n + 1}.

\item On calcule les premiers termes (arrondis au centième) :
\begin{center}
\begin{tabular}{|c|*{9}{c|}}
\hline
$n$ & 0 & 1 & 2 & 3 & 4 & 5 & 6 & 7 & 8 \\ \hline
$u_n$ & 20 & 19,72 & 19,44 & 19,17 & 18,90 & 18,64 & 18,38 & 18,12 & 17,87 \\ \hline
\end{tabular}
\end{center}
On a $u_7 \approx 18{,}12 > 18$ et $u_8 \approx 17{,}87 \leqslant 18$.
Comme la suite est décroissante, $n=8$ est la première valeur pour laquelle $u_n\leqslant 18$.
Le programme renvoie donc $8$ : le chauffage se remet en marche au bout de
\textbf{8 dizaines de minutes}, soit \textbf{80 minutes}.
\end{enumerate}

\end{enumerate}

%==========================================================
\newpage
\exo{Exercice 2 (Métropole 1 2026)}
%==========================================================

$f(x) = a + \dfrac{b\ln(x+1)}{x+1}$ sur $]-1;+\infty[$.

\subsection*{Partie A}

\begin{enumerate}

\item \textbf{Justifions que $a=1$.}

Le point $A(0;1)$ appartient à la courbe, donc $f(0)=1$. Or
\[
f(0) = a + \dfrac{b\ln(1)}{1} = a + 0 = a.
\]
Donc $a = 1$.

\item \textbf{Lectures graphiques.}

\begin{enumerate}
\item $f'(0)$ est le coefficient directeur de la tangente $T_A$ à la courbe en $A(0;1)$.
Sur le graphique, $T_A$ passe par $A(0;1)$ et par le point de coordonnées $(1\,;\,5)$. Donc
\[
f'(0) = \dfrac{5-1}{1-0} = 4.
\]

\item Sur le graphique, la courbe est \emph{concave} au voisinage de $x=1$ : la pente de la tangente
diminue (elle vaut $4$ en $0$ et s'annule au maximum de la courbe, après $x=1$).
Donc $f'$ est décroissante autour de $1$, ce qui signifie $f''(1) < 0$.
\end{enumerate}

\item \textbf{Calcul de $f'$ et valeur de $b$.}

\begin{enumerate}
\item Posons $u(x)=\ln(x+1)$ et $v(x)=x+1$, de dérivées $u'(x)=\dfrac{1}{x+1}$ et $v'(x)=1$.
Pour $x>-1$, avec $f(x)=1+b\,\dfrac{u(x)}{v(x)}$ :
\[
f'(x) = b\times\dfrac{u'(x)v(x)-u(x)v'(x)}{v(x)^2}
= b\times\dfrac{\dfrac{1}{x+1}\times(x+1) - \ln(x+1)\times 1}{(x+1)^2}
= \dfrac{b\bigl(1-\ln(x+1)\bigr)}{(x+1)^2}.
\]

\item D'après la question 2.(a), $f'(0)=4$. Or
\[
f'(0) = \dfrac{b\,(1-\ln 1)}{1^2} = b.
\]
Donc $b = 4$.
\end{enumerate}

\end{enumerate}

\subsection*{Partie B}

On a désormais $f(x) = 1 + \dfrac{4\ln(x+1)}{x+1}$ sur $]-1;+\infty[$.

\begin{enumerate}

\item \textbf{Asymptote.}

On pose $X = x+1$ ; quand $x\to+\infty$, $X\to+\infty$. Par croissances comparées,
$\lim_{X\to+\infty}\dfrac{\ln X}{X}=0$, donc
\[
\lim_{x\to+\infty} \dfrac{4\ln(x+1)}{x+1} = 0
\qquad\text{et}\qquad
\lim_{x\to+\infty} f(x) = 1.
\]
La droite d'équation $y=1$ est donc asymptote horizontale à la courbe de $f$ en $+\infty$.

\item \textbf{Inéquation $1-\ln(x+1)>0$ sur $]-1;+\infty[$.}

\[
1-\ln(x+1)>0 \iff \ln(x+1)<1 \iff x+1<\e \iff x<\e-1,
\]
car la fonction exponentielle est strictement croissante. L'ensemble des solutions est
$\;]-1\,;\,\e-1[$.

\item \textbf{Tableau de variation.}

D'après la partie A, avec $b=4$ : $f'(x)=\dfrac{4\bigl(1-\ln(x+1)\bigr)}{(x+1)^2}$.
Comme $(x+1)^2>0$ sur $]-1;+\infty[$, $f'(x)$ est du signe de $1-\ln(x+1)$ :
$f'>0$ sur $]-1;\e-1[$, $f'(\e-1)=0$, et $f'<0$ sur $]\e-1;+\infty[$.

Extremum (maximum) :
\[
f(\e-1) = 1 + \dfrac{4\ln(\e)}{\e} = 1+\dfrac{4}{\e}.
\]

\begin{center}
\begin{tikzpicture}
\tkzTabInit[lgt=2,espcl=3.5]{$x$/1,$f'(x)$/1,$f(x)$/2}{$-1$,$\e-1$,$+\infty$}
\tkzTabLine{d,+,z,-,}
\tkzTabVar{-/$-\infty$, +/$1+\dfrac{4}{\e}$, -/$1$}
\end{tikzpicture}
\end{center}

(On utilise $\lim_{x\to-1}f(x)=-\infty$ admis, et $\lim_{x\to+\infty}f(x)=1$ vu en question 1.)

\item \textbf{Équation $f(x)=1{,}5$ sur $[2;+\infty[$.}

On a $2>\e-1\approx 1{,}72$, donc $f$ est continue et strictement décroissante sur $[2;+\infty[$.
De plus
\[
f(2) = 1+\dfrac{4\ln 3}{3}\approx 2{,}46 \qquad\text{et}\qquad \lim_{x\to+\infty}f(x)=1,
\]
donc $f$ réalise une bijection de $[2;+\infty[$ sur $]1\,;\,f(2)]$.
Comme $1{,}5\in\,]1\,;\,f(2)]$, d'après le corollaire du théorème des valeurs intermédiaires,
l'équation $f(x)=1{,}5$ admet une unique solution $\alpha$ dans $[2;+\infty[$.

À la calculatrice :
\[
f(25{,}05)\approx 1{,}5006>1{,}5 \qquad\text{et}\qquad f(25{,}15)\approx 1{,}4993<1{,}5.
\]
Comme $f$ est décroissante, $25{,}05<\alpha<25{,}15$, donc $\alpha\approx 25{,}1$ (à $10^{-1}$ près).
\end{enumerate}

%==========================================================
\newpage
\exo{Exercice 3 (Métropole 2 2026)}
%==========================================================

\subsection*{Partie A}

\begin{enumerate}

\item \textbf{Lectures graphiques.}

\begin{enumerate}
\item La courbe $\mathcal{C}_f$ admet au point $C$ d'abscisse $1$ une tangente parallèle à l'axe des abscisses,
donc son coefficient directeur est nul : $f'(1)=0$.

\item Le point $A\left(\tfrac12;0\right)$ appartient à $\mathcal{C}_f$ et c'est le point d'intersection
avec l'axe des abscisses sur $[0;3]$. L'équation $f(x)=0$ a donc pour solution $x=\dfrac12$.
\end{enumerate}

\item \textbf{Détermination de $f'\left(\tfrac12\right)$.}

$f'\left(\tfrac12\right)$ est le coefficient directeur de la tangente $T_A$, qui passe par
$A\left(\tfrac12;0\right)$ et $B(1;\e^2)$ :
\[
f'\!\left(\tfrac12\right)=\dfrac{\e^2-0}{1-\tfrac12}=2\e^2.
\]

\item \textbf{Primitives de $f$.}

Une primitive $F$ de $f$ vérifie $F'=f$. Or, d'après la partie A, $f(x)<0$ sur $\left[0;\tfrac12\right[$
et $f(x)>0$ sur $\left]\tfrac12;+\infty\right[$ (la courbe est sous l'axe avant $A$, au-dessus ensuite).
Donc toute primitive de $f$ est décroissante sur $\left[0;\tfrac12\right]$, puis croissante sur
$\left[\tfrac12;+\infty\right[$ : elle admet un \emph{minimum en $x=\tfrac12$}.
De plus, deux primitives de $f$ diffèrent d'une constante : leurs courbes se déduisent l'une de l'autre
par translation verticale.

Sur la figure, les courbes $\mathcal{C}_1$ et $\mathcal{C}_2$ ont toutes deux leur minimum en
$x=\tfrac12$ et se déduisent l'une de l'autre par une translation verticale, alors que $\mathcal{C}_3$
a son minimum en $x\approx 1$.
\textbf{Les courbes $\mathcal{C}_1$ et $\mathcal{C}_2$ représentent donc des primitives de $f$.}

\emph{Remarque (vérification).} Avec $f(x)=(2x-1)\e^{-2x+3}$, la fonction $F(x)=-x\,\e^{-2x+3}$ est
une primitive de $f$ ; les autres sont de la forme $F+k$.
\end{enumerate}

\subsection*{Partie B}

On admet que $f(x)=(2x-1)\e^{-2x+3}$ sur $[0;+\infty[$.

\begin{enumerate}

\item \begin{enumerate}
\item \textbf{Vérification.} Pour tout $x\geqslant0$ :
\[
f(x)=(2x-1)\e^{-2x+3}=(2x-1)\,\e^{-(2x-1)+2}=\e^{2}\,(2x-1)\,\e^{-(2x-1)}=\e^2\times\dfrac{2x-1}{\e^{2x-1}}.
\]

\item \textbf{Limite en $+\infty$.} On pose $X=2x-1$ ; quand $x\to+\infty$, $X\to+\infty$ et
\[
f(x)=\e^2\times\dfrac{X}{\e^{X}}=\e^2\times\dfrac{1}{\dfrac{\e^X}{X}}.
\]
Par croissances comparées, $\lim_{X\to+\infty}\dfrac{\e^X}{X}=+\infty$, donc
$\lim_{x\to+\infty}f(x)=0$.
\end{enumerate}

\item \begin{enumerate}
\item \textbf{Dérivée.} $f=uv$ avec $u(x)=2x-1$, $u'(x)=2$, $v(x)=\e^{-2x+3}$, $v'(x)=-2\,\e^{-2x+3}$ :
\[
f'(x)=2\,\e^{-2x+3}+(2x-1)\times(-2)\,\e^{-2x+3}
=\bigl(2-4x+2\bigr)\e^{-2x+3}=(-4x+4)\,\e^{-2x+3}.
\]

\item \textbf{Tableau de variation.} Comme $\e^{-2x+3}>0$, $f'(x)$ est du signe de $-4x+4$ :
positif pour $x<1$, nul en $1$, négatif pour $x>1$. De plus
$f(0)=(-1)\e^{3}=-\e^3$, $f(1)=\e^{1}=\e$ et $\lim_{x\to+\infty}f(x)=0$.

\begin{center}
\begin{tikzpicture}
\tkzTabInit[lgt=2,espcl=3.5]{$x$/1,$f'(x)$/1,$f(x)$/2}{$0$,$1$,$+\infty$}
\tkzTabLine{,+,z,-,}
\tkzTabVar{-/$-\e^3$, +/$\e$, -/$0$}
\end{tikzpicture}
\end{center}
\end{enumerate}

\item \textbf{Contrainte du client et détermination de $\alpha$.}

À l'abscisse $\alpha$, le contour est formé du point $(\alpha\,;\,f(\alpha))$ et de son symétrique par
rapport à l'axe des abscisses $(\alpha\,;\,-f(\alpha))$. L'ouverture du logo mesure donc
$f(\alpha)-(-f(\alpha))=2f(\alpha)$. L'unité étant le centimètre, la contrainte « $0{,}3$ cm »
se traduit par $2f(\alpha)=0{,}3$, soit
\[
f(\alpha)=0{,}15.
\]

\emph{Existence et unicité.} Sur $[1;+\infty[$, $f$ est continue et strictement décroissante
(d'après le tableau), avec $f(1)=\e\approx2{,}72$ et $\lim_{x\to+\infty}f(x)=0$.
Donc $f$ réalise une bijection de $[1;+\infty[$ sur $]0\,;\,\e]$. Comme $0{,}15\in\,]0\,;\,\e]$,
l'équation $f(x)=0{,}15$ admet une unique solution $\alpha\in[1;+\infty[$ (corollaire du théorème des
valeurs intermédiaires).

\emph{Valeur approchée.} À la calculatrice :
\[
f(3{,}3)\approx0{,}1530>0{,}15,\qquad f(3{,}35)\approx0{,}1409<0{,}15.
\]
Comme $f$ est décroissante, $3{,}3<\alpha<3{,}35$, donc
\[
\alpha\approx3{,}3\quad\text{(arrondi au dixième).}
\]
\end{enumerate}

\end{document}